\section{分裂半单李代数上的有限维模}

本节中，我们保留 §6 的一般记号。用 P（相应地 Q）记 R 的权格（相应地根权格）。用 \(P_{+}\) （相应地 \(Q_{+}\) ）记 P（相应地 Q）中对由 B 定义的序关系为正的元素组成的集合。用 \(P_{++}\) 记相对于 B 的 R 的支配权组成的集合（第六章 §1 第 10 节）。\(\lambda\) 是 \(h^{*}\) 的一个元素，它属于 P（相应地属于 \(P_{++}\) ）当且仅当所有的 \(\lambda(H_{\alpha})\) （\(\alpha \in B\) ）都是整数（相应地 \(\geq 0\) 的整数）。我们有 \(P_{++} \subset P_{+}\) （第六章 §1 第 6 节）。若 \(w \in W\) ，用 \(\varepsilon(w)\) 记 w 的行列式，它等于 1 或 -1。置 \(\rho = \frac{1}{2} \sum_{\alpha \in R_{+}} \alpha\) 。
% semantic-fragments: legacy-input-seam
\relax{}
